Rapid, equitable, and life-sustaining allocation of emergency relief commodities during acute natural disasters represents one of the most formidable computational and operational challenges faced by municipal governments in disaster-prone developing nations. In uncoordinated crisis response, decentralized request inflation and spatial accessibility bias consistently produce severe logistical failures: accessible urban communities receive surplus aid while isolated, highly vulnerable populations face life-threatening deprivation. Conventional disaster management heuristics either rely on unweighted utilitarian metrics---which induce small-demand paradoxes that starve dense urban hazard centers---or static proportional rationing rules that fail to account for the non-linear, exponential acceleration of human suffering over time.

This study formulates, implements, and benchmarks a centralized, multi-commodity, two-stage Mixed-Integer Linear Programming (MILP) model for the post-disaster Resource Allocation Problem (RAP) across all 182 administrative barangays of Davao City, Philippines. The framework evaluates two life-critical relief commodities---potable water (15~L/person/day) and emergency nutritional subsistence (2,100~kcal/person/day)---parameterized against the 2024 Philippine Statistics Authority (PSA) census and the international Sphere Humanitarian Standards. The model integrates Stage~1 discrete hub activation across six candidate Disaster Risk Reduction and Management Operations Center (DRRMOC) logistics depots subject to an administrative activation cap ($H_{\max} = 4$), and Stage~2 continuous resource dispatch flows. Rather than optimizing purely volumetric throughput, the objective function minimizes aggregate vulnerability-weighted human deprivation costs, modeled via certified piecewise linear outer tangent approximations ($K = 20$ segments) of convex exponential suffering curves ($e^{\lambda_r(1-\phi_{jr})} - 1$). Universal anti-starvation survival floors ($\varepsilon_{\min} = 0.10$) and anti-hoarding demand ceilings ($\phi_{jr} \le 1.0$) are enforced across all barangays.

Empirical evaluation on the baseline disaster instance demonstrates that the proposed Two-Stage MILP achieves certified global optimality ($0.000\%$ relative MIP gap) in \textbf{0.696 seconds} with a memory footprint of \textbf{118.4 MB} on standard local government workstation hardware, exhibiting zero branch-and-bound tree exploration due to root-node LP relaxation convergence. The piecewise linear outer approximation bounds maximum epigraph under-approximation error to $\Delta_{\max} \le 0.005890$ ($< 1.99\%$ for potable water and $< 0.50\%$ for food), preserving over $98.0\%$ to $99.5\%$ mathematical fidelity to the true exponential deprivation penalty. In comparative benchmarking against operational relief policies, the proposed MILP achieves a minimum deprivation cost of $Z^* = \mathbf{2.8283}$, reducing human suffering by \textbf{11.32\%} relative to uniform proportional and egalitarian sharing ($Z = 3.1893$) and by \textbf{1.66\%} relative to a hazard-weighted proportional heuristic ($Z = 2.8761$). 

An objective multi-criteria equity evaluation reveals critical operational trade-offs: while egalitarian baselines achieve higher floor parity ($\phi = 0.3333$) and Nash social welfare, the proposed model deliberately accepts lower floor satisfaction for low-hazard barangays ($56.0\%$ of population at $\phi = 0.10$) to aggressively concentrate scarce marginal supplies into severe-hazard zones ($\bar{\phi}_{\text{High}} = 0.4821$, hazard parity ratio $4.821$). Stress-testing across seven operational calamity scenarios identifies a sharp systemic insolvency boundary at effective supply levels below 10\% (SC-04), functioning as an automated diagnostic trigger for Tier-3 National Escalation. Furthermore, the model demonstrates that identical 60\% deficits produce a 60.8\% higher deprivation penalty for potable water than food ($Z^* = 4.85$ vs. $3.02$), validating physiological urgency prioritization. Spatial analysis confirms that the MILP suppresses geographic accessibility correlation from $r = +1.0000$ to $r = +0.2891$ and prevents decentralized supply over-allocation by construction ($0$ units vs. $2.6$ million wasted units in decentralized simulations). Non-parametric Friedman ($\chi_F^2 = 384.21, p < 0.0001$), post-hoc Nemenyi ($\text{CD} = 0.4518$), and paired Wilcoxon signed-rank tests ($Z \ge 8.42, p < 10^{-6}$) confirm that all allocation shifts are statistically significant across paired barangays. These results provide local emergency authorities with a computationally tractable, mathematically certified, and operationally defensible decision-support framework for city-wide disaster logistics.

\vspace{1.5em}
\noindent\textbf{Keywords:} Humanitarian Logistics, Mixed-Integer Linear Programming, Deprivation Cost Functions, Multi-Commodity Allocation, Distributive Equity, Disaster Operations Research, Davao City.

